\section{Notation and Conventions}
\label{section-notations-conventions}

\noindent
A ring is commutative with $1$. The  maximal ideal of a local ring $A$
is denoted by $\mathfrak{m}_A$. The set of positive integers is denoted
by $\mathbf{N} = \{1, 2, 3, \ldots\}$. If $U$ is an object of a
category $\mathcal{C}$, we denote by $\underline{U}$
the functor
$\Mor_\mathcal{C}(U, -): \mathcal{C} \to \textit{Sets}$, see
Remarks \ref{remarks-cofibered-groupoids} (\ref{item-definition-yoneda}).
Warning: this may conflict with the notation in other chapters where we
sometimes use $\underline{U}$ to denote $h_U(-) = \Mor_\mathcal{C}(-, U)$.

\medskip\noindent
Throughout this chapter $\Lambda$ is a Noetherian ring and
$\Lambda \to k$ is a finite ring map from $\Lambda$ to a field.
The kernel of this map is denoted $\mathfrak m_\Lambda$ and the
image $k' \subset k$. It turns out that $\mathfrak m_\Lambda$ is
a maximal ideal, $k' = \Lambda/\mathfrak m_\Lambda$ is a field, and
the extension $k/k'$ is finite. See discussion surrounding
(\ref{equation-k-prime}).


